\subsection{How fast the losing density vanishes}
Let \(p_-(h)\) be the logarithmic density of \(B_h(t)<0\). On \(I=[0.040,0.041]\), which contains the transition, there is a constant \(C\) such that
\begin{equation}
 |p_-(h)-p_-(h')|\le C|h-h'|^{1/2}\qquad(h,h'\in I).\label{eq:rate}
\end{equation}
Thus the losing density tends to zero at least as fast as \(C\sqrt{h_*-h}\) from below. This estimate addresses the frequency of reversals near the threshold.

To prove it, use the phase variable \(\theta=\gamma_1 t\) on any translated orbit and divide the centered signal by \(2\e^{-h\gamma_1^2}\). The first oscillation is a cosine with amplitude one. The total contribution of the remaining terms to each of its first two derivatives in \(\theta\) has magnitude at most
\[
 \frac{\sum_{\gamma>\gamma_1} m_\gamma\gamma^2\e^{-0.040\gamma^2}}
      {\gamma_1^2\e^{-0.041\gamma_1^2}} < \frac{2}{7}.
\]
Here the first ten certified ordinates and the tail estimate in Proposition~\ref{prop:enclosure} give numerator \(<2\) and denominator \(>7\). For the numerator's tail use
\(32^2(9/4+32^2)\exp(-0.040\cdot32^2)\) since
\(y(9/4+y)\exp(-0.040y)\) decreases for \(y\ge32^2\).
The extra lines in \texttt{certify\_threshold.py} verify both inequalities. The largest of the first two cosine derivatives has modulus at least \(1/\sqrt2\), and the third derivative is uniformly bounded for \(h\in I\). Lemma~\ref{le:occ} with \(K=2\) now gives \(\nu\{z:|f_h(z)|\le s\}\ll\sqrt{s}\) uniformly on \(I\); the amplitude and time normalizations stay in compact positive intervals. Termwise differentiation in \(h\) gives \(\|f_{h'}-f_h\|_\infty\ll|h'-h|\) on \(I\). A sign disagreement forces \(|f_h|\le\|f_{h'}-f_h\|_\infty\), proving \eqref{eq:rate}.
